
Soit \(\underline{M}\) un \(\underline{B}\)-Module gradué libre de type fini. D'après 1.2, on peut trouver un isomorphisme gradué :
\begin{equation}\tag{1.2.2}
  \underline{M}\simeq \bigoplus_{n=n_0}^{n_1} \bigl(\underline{B}(-n)\bigr)^{k_n}.
\end{equation}

On obtient ainsi une décomposition de \(\underline{M}\) en somme directe de \(\underline{B}\)-Modules gradués libres de type fini et homogènes, soit :
\[
  \underline{M}\simeq \bigoplus_{n=n_0}^{n_1}\underline{M}^{(n)},
\]
où \(\underline{M}^{(n)}\) est libre homogène de degré \(n\) et de rang \(k_n\).

\textbf{Proposition 1.2.3.} \uline{Les entiers} \(k_n\) \uline{intervenant dans la décomposition} (1.2.2) \uline{ne dépendent pas de la décomposition choisie}.

Supposons données deux décompositions, donc un isomorphisme gradué :
\[
  \bigoplus_{n=n_0}^{n_1}\bigl(\underline{B}(-n)\bigr)^{k_n}\xrightarrow{\sim}
  \bigoplus_{n=n'_0}^{n'_1}\bigl(\underline{B}(-n)\bigr)^{k'_n}.
\]
La démonstration se fait par récurrence sur le nombre de facteurs directs libres homogènes non nuls intervenant dans la décomposition de gauche. Si c'est \(0\), alors \(\underline{M}=0\), et la proposition est évidente. Supposons la démontrée pour \(n_1-n_0=m\), et montrons la pour \(n_1-n_0=m+1\). L'isomorphisme donné, étant compatible aux graduations, donne un isomorphisme sur les composantes homogènes des deux membres, et en particulier un isomorphisme des composantes homogènes de degré \(<n_0\), ce qui montre que pour \(n<n_0\), on a \(k'_n=0\), et sur les composantes homogènes de degré \(n_0\), ce qui montre qu'elles ont même rang, et par suite que \(k_{n_0}=k'_{n_0}\). L'isomorphisme donné induit donc un isomorphisme des facteurs directs libres homogènes de plus bas degré des deux décompositions et passe donc au quotient par ces facteurs directs; on applique alors l'hypothèse de récurrence au Module quotient, d'où le résultat.

